\scanpage{83}


\deleted{Supposons} Nous cherchons maintenant la condition pour que

a) $E_1\otimes E_2\longrightarrow E_1*E_1$ surjectif

b) Son noyau $N$ (qui contient les $\overline{\xi}-\xi$, $\xi\in E_1\otimes\overline{E_2}$) est engendré par \deleted{les} $\operatorname{Im}(\sigma_1\otimes\sigma_2-\mathrm{id})$, i.e. par les $\overline{x_1}\otimes\overline{x_2}-x_1\otimes x_2$ $(x_1\deleted{\unsure{,}}\in\Gamma E_1,\ x_2\in\Gamma E_2)\deleted{\unsure{x}}$, \deleted{\unsure{ou d'un|condition}} en d'autres termes on veut

\[ E_1\otimes E_2\mathbin{\unsure{-|\to}}E_1*E_2\quad\text{induit}\quad {E_1\otimes E_2}_{/\sigma_1\otimes\sigma_2}\simeq E_1*\unsure{I_2|E_2} \]

\deleted{\unsure{a}} Pour la condition a), elle équivaut à la surjectivité de

\[ L_1\otimes E_2\underset{L_1\otimes L_2}{\oplus}E_1\otimes L_2\longrightarrow L_1\otimes L_2,\]

or son conoyau est can. iso. à

\[ \operatorname{Coker}\pi_1\otimes\operatorname{Coker}\pi_2\qquad \pi_i:E_i\longrightarrow L_i \]

Si $L_1,L_2$ sont de présentation finie, alors la nullité (par Nakayama) se vérifie point par point, donc il est non seulement suffisant \inserted{(\unsure{tjs|après} OK)}, mais aussi néc., que \inserted{au voisinage de chaque point de $X$} $\pi_1$ \uuline{OU} $\pi_2$ soient épi. Si on a des scindages $x_1,x_2$,

d'où \quad $b_1\in\Gamma L_1,\ b_2\in\Gamma L_2$

alors

\[ \operatorname{Im}\deleted{\unsure{\pi_i:\,-}}(\pi_i:E_i\mathbin{\unsure{-|\to}}L_i)=\mathcal{O}_X\cdot b_i+2L_i. \]

donc la condition est la suivante :
% slanted

\margin{NB Cette \unsure{traduction|condition} N'utilise que la \unsure{condition|réduction} \unsure{sur|pour} $p=2$}

\textcircled{a$'$} \inserted{En tout} \deleted{Aux} point\deleted{s} de $X_0=V(\unsure{p|\gamma|2|\chi}\,.1)$, \deleted{\unsure{ou bien $b_1$ ou|ou $b_1$ ou $b_2$, ou}} \inserted{\unsure{Resp. ou|Resp. aux}} $T_{10}$, ou bien $T_{20}$, est un \deleted{\unsure{inv|base}} \inserted{gén.} de $L_{10}$ resp $L_{20}$.
% archive sheet number: 82
